% PETUNJUK AKHIR final-pipeline, GEMASTIK XIX Penambangan Data (11 November 2026).
% Bangun: pdflatex PETUNJUK_AKHIR.tex (dua kali).
\documentclass[11pt,a4paper]{article}
\usepackage[left=2.5cm,right=2.5cm,top=2.3cm,bottom=2.4cm]{geometry}
\usepackage[T1]{fontenc}
\usepackage[utf8]{inputenc}
\usepackage[indonesian]{babel}
\usepackage{lmodern}
\usepackage{amssymb,booktabs,tabularx,array,enumitem,titlesec,float,xcolor}
\usepackage[most]{tcolorbox}
\usepackage{tikz}
\usetikzlibrary{arrows.meta,positioning,shapes.multipart,calc}
\usepackage[hidelinks]{hyperref}
\setlength{\parindent}{0pt}
\setlength{\parskip}{5pt}
\setlength{\emergencystretch}{2em}
\titleformat{\section}{\Large\bfseries}{\thesection}{0.7em}{}
\titleformat{\subsection}{\large\bfseries}{\thesubsection}{0.6em}{}
\titlespacing*{\section}{0pt}{14pt}{6pt}
\titlespacing*{\subsection}{0pt}{10pt}{3pt}
\setlist{leftmargin=*,itemsep=1pt,parsep=0pt,topsep=2pt}
\renewcommand{\arraystretch}{1.2}
\newcolumntype{Y}{>{\raggedright\arraybackslash}X}

% ---------- warna ----------
\definecolor{biru}{HTML}{1F5FAE}
\definecolor{hijau}{HTML}{1A9B6C}
\definecolor{jingga}{HTML}{E8663D}
\definecolor{kuning}{HTML}{C8860A}
\definecolor{ungu}{HTML}{7A4FB5}
\definecolor{merah}{HTML}{C0392B}
\definecolor{abu}{HTML}{6B7280}

% ---------- alat tulis ----------
\newcommand{\f}[1]{\texttt{\detokenize{#1}}}                 % nama file / kode
\newcommand{\ket}[1]{{\color{black!55}#1}}                    % keterangan abu-abu
\newcommand{\cek}{$\square$\enspace}
\newcommand{\pil}[2]{\tikz[baseline=(p.base)]\node[fill=#1,text=white,font=\bfseries\scriptsize,
  rounded corners=2pt,inner xsep=4pt,inner ysep=2pt](p){#2};}

% kotak gambar bernomor, pita judul abu-abu
\newtcolorbox[auto counter]{gambar}[2][]{enhanced,colback=white,colframe=black!18,boxrule=0.5pt,
  arc=2pt,colbacktitle=black!6,coltitle=black,fonttitle=\bfseries\small,
  title={Gambar \thetcbcounter. #2},left=6pt,right=6pt,top=6pt,bottom=6pt,
  before skip=8pt,after skip=10pt,halign=center,#1}
% kartu notebook: #1 warna, #2 nama, #3 keterangan kanan
\newtcolorbox{nb}[3]{enhanced,colback=white,colframe=#1,boxrule=0.5pt,leftrule=4pt,arc=2pt,
  colbacktitle=#1!10,coltitle=#1!75!black,fonttitle=\bfseries,
  title={#2\hfill{\normalfont\small #3}},left=6pt,right=6pt,top=3pt,bottom=3pt,
  before skip=7pt,after skip=7pt,fontupper=\small}
\newenvironment{isi}{\tabularx{\linewidth}{@{}>{\bfseries}p{1.45cm}Y@{}}}{\endtabularx}
% kotak peringatan / catatan
\newtcolorbox{awas}[1][merah]{enhanced,colback=#1!5,colframe=#1,boxrule=0pt,leftrule=3pt,arc=0pt,
  left=7pt,right=7pt,top=4pt,bottom=4pt,before skip=7pt,after skip=7pt}

\tikzset{
  ar/.style={-{Stealth[length=2mm]},black!60,line width=0.7pt},
  kartu/.style 2 args={rectangle split,rectangle split parts=2,draw=#1,rounded corners=2pt,
    rectangle split part fill={#1,white},text width=#2,align=left,font=\scriptsize,inner sep=3pt},
  kotak/.style={draw=black!45,rounded corners=2pt,align=center,font=\footnotesize,inner sep=4pt},
  hasil/.style={draw=#1,fill=#1!12,rounded corners=2pt,align=left,font=\footnotesize,inner sep=4pt},
}
\newcommand{\jdl}[1]{\textcolor{white}{\textbf{#1}}}

% ---------- angka gladi (diukur 10 November 2026 di laptop, data demo SmSA) ----------
\newcommand{\wEda}{23 s}\newcommand{\wPra}{60 s}\newcommand{\wBase}{24 s, F1 OOF 0,860}
\newcommand{\wFt}{31 s per epoch per fold; 1 epoch F1 OOF 0,927}
\newcommand{\wEns}{5 s}\newcommand{\wTopik}{5 menit}\newcommand{\wBer}{80 s}\newcommand{\wLap}{7 s}

\begin{document}

% ================= kepala =================
\begin{tcolorbox}[enhanced,colback=biru,colframe=biru,arc=3pt,left=12pt,right=12pt,top=10pt,bottom=10pt]
  {\color{white}{\LARGE\bfseries Petunjuk Akhir: final-pipeline}\\[3pt]
  {\large GEMASTIK XIX Penambangan Data \textbullet\ Rabu, 11 November 2026 \textbullet\ 13.00--18.00}\\[6pt]
  \small Tebakan soal: \textbf{klasifikasi teks} (skor di Kaggle), single label atau multilabel, dan
  \textbf{topic modeling} (dipresentasikan). Langkah bernomor; nilai soal diisi di sel Konfigurasi tiap notebook.}
\end{tcolorbox}

\begin{gambar}{Lima jam dalam satu garis (rencana, bukan aturan panitia)}
\begin{tikzpicture}[x=0.05cm,y=1cm]
  % 1 menit = 0,05 cm
  \foreach \a/\b/\c/\n in {0/20/biru/1,20/45/hijau/2,45/75/kuning/3,75/180/jingga/4,
                           180/210/ungu/5+7,210/240/abu/8,240/300/merah/submit}
    {\fill[\c,rounded corners=1.5pt] (\a+0.6,0) rectangle (\b-0.6,0.75);
     \node[text=white,font=\bfseries\small] at ({(\a+\b)/2},0.375) {\n};}
  \fill[hijau!70!black,rounded corners=1.5pt] (85,-0.62) rectangle (170,-0.12);
  \node[text=white,font=\bfseries\scriptsize] at (127.5,-0.37) {6 topik, selagi GPU melatih};
  \foreach \m/\j in {0/13.00,60/14.00,120/15.00,180/16.00,240/17.00,300/18.00}
    {\draw[black!50] (\m,0.85)--(\m,1.0); \node[font=\scriptsize,above] at (\m,0.95) {\j};}
  \draw[black!55,line width=0.6pt] (0,1.55)--(238,1.55) node[midway,above,font=\scriptsize]
    {tanpa internet, tanpa leaderboard: putusan dari CV lokal};
  \draw[merah,line width=0.9pt] (242,1.55)--(300,1.55) node[midway,above,font=\scriptsize,text=merah]
    {internet, 3 submission};
\end{tikzpicture}

\vspace{2pt}
{\scriptsize\raggedright
\pil{biru}{1} soal, isi NILAI SOAL, EDA\quad
\pil{hijau}{2} praproses\quad
\pil{kuning}{3} baseline + submission pertama di disk\quad
\pil{jingga}{4} fine-tune (1--2 model)\quad
\pil{ungu}{5+7} ensemble, laporan\quad
\pil{abu}{8} pilih 3 file, ekspor\par}
\end{gambar}

% ================= 1 =================
\section{Peta pipeline}

Folder \f{final-pipeline} berisi salinan tujuh notebook repo yang sudah teruji, diberi nomor urut.
Kodenya sama dengan \f{contoh/text/}; yang berbeda hanya path dan sel Konfigurasi, yang diawali
blok \textbf{NILAI SOAL}. Tidak ada \f{konfig.py}: tiap notebook berdiri sendiri. Jalankan menurut nomor.

\begin{gambar}{Alur file dari soal sampai bahan presentasi}
\begin{tikzpicture}
  \node[fill=black!80,text=white,rounded corners=2pt,font=\footnotesize,minimum width=15.5cm,
        minimum height=6mm,anchor=north west] (k) at (0,1.5)
    {\textbf{NILAI SOAL}\quad awal sel Konfigurasi tiap notebook: file soal, kolom, folder hasil (isi sama di semua)};
  \node[kartu={biru}{2.45cm},anchor=north west] (n1) at (0,0)
    {\jdl{\f{1_eda}}\nodepart{two}kenali data, cek bocor, train vs test\\ \ket{temuan.md}};
  \node[kartu={hijau}{2.45cm},anchor=north west] (n2) at (3.2,0)
    {\jdl{\f{2_praproses}}\nodepart{two}periksa kotoran, bersih ringan + berat\\ \ket{\f{train_bersih.csv}}};
  \node[kartu={kuning}{2.45cm},anchor=north west] (n3) at (6.4,0)
    {\jdl{\f{3_baseline}}\nodepart{two}TF-IDF + SVM, 5 fold, validator\\ \ket{oof + submission}};
  \node[kartu={jingga}{2.45cm},anchor=north west] (n4) at (9.6,0)
    {\jdl{\f{4_finetune}}\nodepart{two}IndoBERT dkk., 5 fold, GPU\\ \ket{oof + submission}};
  \node[kartu={ungu}{2.45cm},anchor=north west] (n5) at (12.8,0)
    {\jdl{\f{5_ensemble}}\nodepart{two}gabung OOF, skor jujur\\ \ket{submission akhir}};
  \node[kartu={hijau!70!black}{2.45cm},anchor=north west] (n6) at (3.2,-2.9)
    {\jdl{\f{6_topik}}\nodepart{two}NMF / LDA / embedding, pilih jumlah topik\\ \ket{\f{fakta_topik.md}}};
  \node[kartu={hijau!70!black}{2.45cm},anchor=north west] (n6b) at (6.4,-2.9)
    {\jdl{\f{6b_bertopic}}\nodepart{two}BERTopic (HDBSCAN), pembanding\\ \ket{\f{info_topik.csv}}};
  \node[kartu={ungu}{2.45cm},anchor=north west] (n7) at (9.6,-2.9)
    {\jdl{\f{7_laporan}}\nodepart{two}per kelas, confusion matrix, ECE\\ \ket{\f{fakta_model.md}}};
  \node[kartu={abu}{2.45cm},anchor=north west] (n8) at (12.8,-2.9)
    {\jdl{\f{8_ekspor.bat}}\nodepart{two}saring yang boleh keluar\\ \ket{\f{bawa_pulang/}}};
  \foreach \a/\b in {n1/n2,n2/n3,n3/n4,n4/n5} \draw[ar] (\a.east)--(\b.west);
  \foreach \n in {n1,n2,n3,n4,n5} \draw[ar,black!35] (\n.north |- k.south)--(\n.north);
  \draw[ar] (n2.south)--(n6.north);
  \draw[ar] (n4.south)--(n7.north) node[midway,right,font=\tiny] {OOF};
  \draw[ar] (n5.south)--(n8.north);
  \draw[ar,densely dashed] (n6.east)--(n6b.west);
  \draw[ar] (n7.east)--(n8.west);
  \draw[ar] (n6.south)-- ++(0,-0.35) -| (n8.south);
\end{tikzpicture}
\end{gambar}

\begin{gambar}{Susunan folder yang wajib (path di notebook relatif)}
\begin{minipage}{0.56\linewidth}\small\ttfamily\raggedright
C:/GEMASTIK/\\
\quad not-a-super-secret-weapon/\quad\ket{\rmfamily repo}\\
\quad\quad \textbf{final-pipeline/}\quad\ket{\rmfamily folder ini}\\
\quad\quad\quad 1\_eda.ipynb \dots\ 7\_laporan\_model.ipynb\\
\quad\quad\quad 3m\_, 4m\_, 5m\_ \dots\ \ket{\rmfamily multilabel}\\
\quad\quad\quad 8\_ekspor.bat, jupyter.bat\\
\quad\quad\quad demo/, kamus/\\
\quad pretrained/\quad\ket{\rmfamily model (\f{download_model.bat})}\\
\quad .venv/\quad\ket{\rmfamily dibuat setup.bat}\\
\quad outputs/final/\quad\ket{\rmfamily semua hasil}
\end{minipage}\hfill
\begin{minipage}{0.4\linewidth}\small\raggedright
\textbf{Buka:} dobel klik \f{jupyter.bat} di folder ini.\par\smallskip
\textbf{Tidak butuh} \f{data_latihan}: data demo dan kamus alay sudah ikut di sini.\par\smallskip
Folder ini \textbf{ada di dalam repo}, jadi ikut \f{git clone} / \f{git pull}.
\end{minipage}
\end{gambar}

% ================= 2 =================
\section{Yang diubah saat soal dibuka}

Blok \textbf{NILAI SOAL} di awal sel Konfigurasi (sel kode pertama sesudah judul ``Konfigurasi'')
di \textbf{setiap} notebook yang dijalankan. Tidak ada \f{konfig.py}; tiap notebook menyimpan nilainya
sendiri, jadi \textbf{isi nilai yang sama di semua notebook}, terutama \f{HASIL}: notebook 3--7 membaca
\f{<HASIL>/praproses/train_bersih.csv} buatan \f{2_praproses}, dan \f{HASIL} yang beda membuat file itu
tidak ketemu. Ubah baris yang ada saja; jangan menyalin blok utuh dari notebook lain, karena tiap
notebook hanya memuat variabel yang dipakainya (Tabel di bawah). Sesudah mengubah: \emph{Restart
Kernel and Run All}. Di bawah blok ada bagian ``jalur final-pipeline'' (path dari \f{HASIL}): tidak
perlu diubah.

\begin{tcolorbox}[enhanced,colback=black!3,colframe=black!25,boxrule=0.5pt,arc=2pt,
  left=8pt,right=8pt,top=5pt,bottom=5pt,fontupper=\small\ttfamily]
\textcolor{biru}{TRAIN\_MENTAH} = "C:/GEMASTIK/soal/train.csv"\hfill\pil{biru}{wajib}\\
\textcolor{biru}{TEST\_MENTAH}\ \ = "C:/GEMASTIK/soal/test.csv"\hfill\pil{biru}{wajib}\\
\textcolor{biru}{SAMPLE\_PANITIA} = "C:/GEMASTIK/soal/sample\_submission.csv"\hfill\pil{biru}{wajib}\\[3pt]
\textcolor{hijau}{ID\_COL} = "id"\qquad \textcolor{hijau}{TEXT\_COLS} = ["text"]\qquad
\textcolor{hijau}{TARGET} = "label"\hfill\pil{hijau}{wajib}\\
\textcolor{hijau}{SUBMIT\_COL} = TARGET\hfill\pil{abu}{bila nama di sample beda}\\[3pt]
\textcolor{jingga}{HASIL} = "../../outputs/lomba"\hfill\pil{jingga}{ganti dari "final"}\\
\textcolor{abu}{LANGUAGE}, \textcolor{abu}{DATE\_COL}, \textcolor{abu}{GROUP\_COL},
\textcolor{abu}{TFIDF\_BERAT}, \textcolor{abu}{TOPIK\_BERAT}\hfill\pil{abu}{bila perlu}
\end{tcolorbox}

{\small\renewcommand{\arraystretch}{1.15}
\noindent\begin{tabularx}{\linewidth}{@{}>{\raggedright\arraybackslash}p{4.1cm}X@{}}
\toprule
\textbf{notebook} & \textbf{baris di blok NILAI SOAL} \\ \midrule
\f{1_eda} & \f{TRAIN_MENTAH TEST_MENTAH ID_COL TEXT_COLS TARGET SUBMIT_COL LABEL_COLS DATE_COL LANGUAGE HASIL SCORING} \\
\f{2_praproses} & \f{TRAIN_MENTAH TEST_MENTAH ID_COL TEXT_COLS TARGET SUBMIT_COL LABEL_COLS HASIL SCORING} \\
\f{3_baseline} & \f{SAMPLE_PANITIA ID_COL TEXT_COLS TARGET SUBMIT_COL LABEL_COLS GROUP_COL DATE_COL HASIL TFIDF_BERAT} \\
\f{4_finetune} & \f{SAMPLE_PANITIA ID_COL TEXT_COLS TARGET SUBMIT_COL LABEL_COLS GROUP_COL HASIL} \\
\f{5_ensemble} & \f{SAMPLE_PANITIA ID_COL TEXT_COLS TARGET SUBMIT_COL LABEL_COLS HASIL} \\
\f{6_topik} & \f{ID_COL TEXT_COLS TARGET SUBMIT_COL LABEL_COLS DATE_COL LANGUAGE HASIL TOPIK_BERAT} \\
\f{6b_topik_bertopic} & \f{TEXT_COLS TARGET SUBMIT_COL LABEL_COLS LANGUAGE HASIL} \\
\f{7_laporan_model} & \f{TEXT_COLS TARGET SUBMIT_COL LABEL_COLS HASIL} \\
\f{3m_baseline_multilabel} & \f{SAMPLE_PANITIA ID_COL TEXT_COLS TARGET SUBMIT_COL LABEL_COLS HASIL TFIDF_BERAT} \\
\f{4m_finetune_multilabel} & \f{SAMPLE_PANITIA ID_COL TEXT_COLS TARGET SUBMIT_COL LABEL_COLS HASIL} \\
\f{5m_ensemble_multilabel} & \f{SAMPLE_PANITIA ID_COL TARGET SUBMIT_COL LABEL_COLS HASIL} \\
\f{9_figur_ppt.py} & \f{HASIL} (baris atas file, bagian ``Konfigurasi'') \\
\f{9b_slide.py} & \f{HASIL ID_COL TEXT_COLS TARGET} (baris atas file, bagian ``Konfigurasi'') \\
\f{8_ekspor.bat} & baris \f{set "HASIL=..."} \\
\bottomrule
\end{tabularx}}

\begin{awas}
\textbf{Tiga hal di luar blok NILAI SOAL.}
\begin{enumerate}
  \item \textbf{Metrik.} Bawaannya F1 makro. Kalau soal memakai metrik lain, pindahkan \texttt{\#} di
  baris \f{return} sel ``Metrik'' pada \f{3_baseline}, \f{4_finetune}, dan \f{5_ensemble}
  (ketiganya harus sama).
  \item \textbf{Pemisah dan encoding CSV.} Kalau \f{pd.read_csv} gagal atau kolom jadi satu,
  pakai baris ber-\texttt{\#} di sel ``Baca data'' \f{2_praproses} (\f{sep=";"},
  \f{encoding="latin-1"}). Notebook 3--7 membaca \f{train_bersih.csv}, jadi cukup dibetulkan di
  notebook 1 dan 2.
  \item \textbf{Pola awalan/akhiran.} \f{STRIP_PATTERNS} di \f{2_praproses} masih pola demo.
  Isi dari tabel ``Awalan dan akhiran berulang'' di notebook itu, atau kosongkan (\f{[]}).
\end{enumerate}
\end{awas}

% ================= 3 =================
\section{Langkah demi langkah}

Waktu di pojok kanan diukur pada gladi 10 November 2026 di laptop (RTX 5060 8 GB) dengan data
demo SmSA: 11.000 teks latih, 3 kelas. Data lomba yang lebih besar atau lebih panjang akan lebih
lama.

\begin{nb}{biru}{\f{1_eda}}{demo \wEda}
\begin{isi}
Tujuan & Tahu bentuk data sebelum memilih apa pun. Tidak mengubah file.\\
Lihat & Sebaran kelas, panjang teks, teks kembar dan kembar beda label, kekuatan tiap kolom
(tanda bocor), AUC adversarial train vs test.\\
Luaran & \f{<HASIL>/eda/temuan.md} dan PNG. \f{AMAN_NDA = True}: tidak mencetak teks asli.
Soal multilabel: bagian 7f (lihat bagian ``Kalau soalnya multilabel'').
\end{isi}
\end{nb}

\begin{nb}{hijau}{\f{2_praproses}}{demo \wPra}
\begin{isi}
Tujuan & Membuat \f{train_bersih.csv} dan \f{test_bersih.csv} yang dibaca semua notebook
sesudahnya.\\
Ubah & \f{STRIP_PATTERNS}; saklar \f{DROP_*} di bagian 7 (teks kembar beda label).\\
Lihat & Tabel ``Ukur'': langkah mana yang menaikkan skor. Tiga baris \f{CEK} di sel terakhir
harus \f{True}.\\
Awas & Jumlah baris train bisa berkurang (kosong, kembar). Jangan jalankan ulang notebook ini
sesudah ada OOF, kecuali semua model dilatih ulang.
\end{isi}
\end{nb}

\begin{nb}{kuning}{\f{3_baseline}}{demo \wBase}
\begin{isi}
Tujuan & Submission pertama yang sah dan OOF pertama. Target: selesai sebelum 14.15.\\
Ubah & \f{MODEL_NAME} tiap variasi (mis. \f{tfidf_logreg}), lalu pindahkan \texttt{\#} di
\f{CLASSIFIER}. Nama sama berarti OOF lama tertimpa.\\
Lihat & Skor OOF dan baris terakhir validator: \textbf{\f{MASALAH: 0}}. Tulisan
\f{PERINGATAN: SAMPLE_CSV masih file demo} berarti \f{SAMPLE_PANITIA} belum diisi.\\
Luaran & \f{baseline/submission.csv}, \f{oof/oof_<nama>.npy}, \f{oof/percobaan.csv}.
\end{isi}
\end{nb}

\begin{nb}{jingga}{\f{4_finetune}}{demo \wFt}
\begin{isi}
Tujuan & Skor tertinggi. Zero-shot dan few-shot dilarang; fine-tune sendiri seperti ini boleh.\\
Ubah & \f{PRETRAINED} + \f{MODEL_NAME} (Gambar~\ref{g:model}); \f{MAX_LENGTH} dari tabel
persentil token (ambil sekitar p95).\\
Uji dulu & \f{FOLDS_TO_RUN = [0]} untuk mengukur waktu satu fold (OOF tidak disimpan), lalu
kembalikan ke \f{None}. Waktu penuh kira-kira 5 kali satu fold. Gladi memakai 1 epoch (3 menit); bawaan 3 epoch, jadi kira-kira 8 menit di data demo.\\
Awas & GPU hanya untuk satu notebook fine-tune pada satu waktu. Model \emph{large}:
\f{BATCH_SIZE = 8}, \f{ACCUM_STEPS = 2}, \f{GRAD_CHECKPOINT = True}.
\end{isi}
\end{nb}

\begin{gambar}[label=g:model]{Memilih \texttt{PRETRAINED} di \texttt{4\_finetune}}
\begin{tikzpicture}
  \node[kotak,text width=2.1cm] (q) at (0,0) {\textbf{Teksnya seperti apa?}\\ \ket{lihat \f{1_eda}}};
  \node[kotak,text width=4.6cm,align=left] (a) at (4.6,1.8) {\textbf{Medsos}: alay, emoji, mention};
  \node[kotak,text width=4.6cm,align=left] (b) at (4.6,0.6) {\textbf{Formal}: berita, ulasan rapi};
  \node[kotak,text width=4.6cm,align=left] (c) at (4.6,-0.6) {\textbf{Campur}: daerah, Inggris};
  \node[kotak,text width=4.6cm,align=left] (d) at (4.6,-1.8) {\textbf{Waktu mepet} atau data sangat besar};
  \node[hasil=jingga,text width=5.3cm] (ra) at (11.2,1.8) {\f{indobertweet} + nyalakan \f{prep_tweet}};
  \node[hasil=biru,text width=5.3cm] (rb) at (11.2,0.6) {\f{indobert}, lalu \f{indobert_large}};
  \node[hasil=hijau,text width=5.3cm] (rc) at (11.2,-0.6) {\f{xlmr_base} atau \f{nusabert}};
  \node[hasil=abu,text width=5.3cm] (rd) at (11.2,-1.8) {\f{distilbert_id} atau \f{minilm}};
  \foreach \x/\y in {a/ra,b/rb,c/rc,d/rd} {\draw[ar] (q.east)--(\x.west); \draw[ar] (\x.east)--(\y.west);}
\end{tikzpicture}

\vspace{2pt}\parbox{\linewidth}{\centering\scriptsize\ket{Dua model yang berbeda (mis.
\f{indobert} dan \f{xlmr_base}) lebih berguna untuk ensemble daripada satu model dengan dua seed.}}
\end{gambar}

\begin{nb}{ungu}{\f{5_ensemble}}{demo \wEns}
\begin{isi}
Tujuan & Menggabung semua OOF di \f{oof/} dan menilai dengan skor jujur (bobot dicari di lipatan
lain).\\
Lihat & Sel ``Putusan''. Gabungan hanya dipakai bila naik sekurangnya \f{MIN_GAIN} (0,002);
kalau tidak, notebook memilih model tunggal terbaik. Ikuti putusannya.\\
Luaran & \f{ensemble/submission.csv}, \f{ablation.csv} (bahan slide).
\end{isi}
\end{nb}

\begin{nb}{hijau!70!black}{\f{6_topik} dan \f{6b_topik_bertopic}}{demo \wTopik\ dan \wBer}
\begin{isi}
Tujuan & Topic modeling untuk presentasi. \f{6_topik} membandingkan NMF, LDA, dan
embedding+KMeans pada beberapa jumlah topik dengan koherensi NPMI, keragaman, dan kestabilan,
lalu memilih sendiri.\\
Masukan & \f{<HASIL>/praproses/train_bersih.csv} dan \f{test_bersih.csv} (dari \f{2_praproses};
belum ada = jalankan \f{2_praproses} dulu). Dua kolom dipakai: \f{text} (bersih ringan) untuk
embedding MiniLM \f{../../pretrained/minilm}, dan \f{text_tfidf} (bersih berat) untuk kata topik.
Nama kolom mengikuti \f{TEXT_COLS} di blok NILAI SOAL \f{6_topik.ipynb}.\\
Kata topik & Saklar \f{TOPIK_BERAT} di blok NILAI SOAL \f{6_topik.ipynb}. \f{True} (bawaan): kata dari kolom
\f{text_tfidf}, jadi slang sudah baku, imbuhan dibuang (stemming), dan ``gak''/``tidak'' atau
``makanannya''/``makan'' jadi satu kata. Kata jadi bentuk dasar (``layan'', bukan ``pelayanan'').
\f{False}: kata dari kolom \f{text}, hanya stopword dibuang; pakai untuk teks baku atau Inggris,
atau bila stemming merusak nama. Isi kamus slang: \f{kamus/kamus_alay_salsabila.csv} dan
\f{kamus/kamus_slang.csv}. Urutan langkah berat: \f{HEAVY_STEPS} di sel bagian 4
\f{2_praproses.ipynb}; diubah = jalankan ulang \f{2_praproses} lalu \f{6_topik}.\\
Ubah & Di sel ``Konfigurasi'' \f{6_topik.ipynb}: \f{EXTRA_STOPWORDS} (kata umum yang mengotori
topik, mis.\ \f{"sih"}), \f{TOPIC_COUNTS} bila jumlah topik yang wajar jauh dari 4--15,
\f{METHODS}. Lalu di sel ``Beri nama topik'': \f{TOPIC_NAMES} sesudah membaca kata teratas
(tulis bentuk wajar, mis.\ ``pelayanan''), jalankan ulang dari sel itu ke bawah.\\
Kapan & Selagi \f{4_finetune} berjalan. NMF dan LDA memakai CPU. Metode \f{"embedding"} dan
\f{6b} memakai MiniLM di GPU; bila VRAM penuh, hapus \f{"embedding"} dari \f{METHODS} dulu.\\
Luaran & Folder \f{<HASIL>/topik/}: \f{fakta_topik.md} (kalimat untuk slide), \f{01_}..\f{09_*.png}
+ \f{.csv} (pemilihan K, ukuran, kata teratas, awan kata, kemiripan, topik $\times$ kelas,
train vs test, tren, peta 2D), \f{info_topik.csv}. \f{topik_dokumen.csv} (topik per baris)
tetap di komputer lomba.\\
\f{6b} & Hanya pembanding, membaca \f{train_bersih.csv} kolom \f{text} (tanpa \f{TOPIK_BERAT}).
BERTopic menentukan jumlah topik sendiri dan menyisakan topik $-1$. Luaran: \f{<HASIL>/bertopic/}.
\end{isi}
\end{nb}

\begin{nb}{ungu}{\f{7_laporan_model}}{demo \wLap}
\begin{isi}
Tujuan & Bahan slide hasil klasifikasi dari OOF, tanpa GPU: skor tiap model, laporan per kelas,
confusion matrix, pasangan kelas tertukar, keandalan (ECE).\\
Luaran & \f{laporan/fakta_model.md} dan 10 grafik. Jalankan sesudah model terakhir selesai.
\end{isi}
\end{nb}

\begin{nb}{abu}{\f{8_ekspor.bat}}{beberapa detik}
\begin{isi}
Tujuan & Menyalin yang boleh keluar ke \f{outputs/bawa_pulang/} dan menolak sisanya.\\
Ubah & Baris \f{set "HASIL=..."} di \f{8_ekspor.bat} bila \f{HASIL} di notebook diganti (isi sama).\\
Lihat & \f{MANIFEST.md}: daftar yang disalin, yang ditolak, dan alasannya.
\end{isi}
\end{nb}

% ================= multilabel =================
\section{Kalau soalnya multilabel}

Multilabel berarti satu teks boleh punya beberapa label sekaligus, dan target berupa beberapa
kolom 0/1. Notebook 3, 4, 5, dan 7 \textbf{tidak} bisa dipakai untuk itu. Gantinya tiga notebook
berakhiran \f{m}; EDA, praproses, dan topic modeling tetap notebook yang sama.

\begin{gambar}{Jalur single label dan jalur multilabel}
\begin{tikzpicture}
  \node[kartu={biru}{2.2cm},anchor=north west] (a1) at (0,0) {\jdl{\f{1_eda}}\nodepart{two}sama};
  \node[kartu={hijau}{2.2cm},anchor=north west] (a2) at (2.9,0) {\jdl{\f{2_praproses}}\nodepart{two}sama};
  \node[kartu={abu}{2.6cm},anchor=north west] (s3) at (6.0,0.95)
    {\jdl{\f{3_baseline}}\nodepart{two}satu kelas per teks};
  \node[kartu={abu}{2.6cm},anchor=north west] (s4) at (9.3,0.95)
    {\jdl{\f{4_finetune}}\nodepart{two}softmax};
  \node[kartu={abu}{2.6cm},anchor=north west] (s5) at (12.6,0.95)
    {\jdl{\f{5_ensemble} + \f{7}}\nodepart{two}gabung, laporan};
  \node[kartu={kuning}{2.6cm},anchor=north west] (m3) at (6.0,-0.95)
    {\jdl{\f{3m_baseline}}\nodepart{two}TF-IDF, satu model per label};
  \node[kartu={jingga}{2.6cm},anchor=north west] (m4) at (9.3,-0.95)
    {\jdl{\f{4m_finetune}}\nodepart{two}sigmoid + BCE, ambang per label};
  \node[kartu={ungu}{2.6cm},anchor=north west] (m5) at (12.6,-0.95)
    {\jdl{\f{5m_ensemble}}\nodepart{two}gabung + laporan per label};
  \node[kartu={hijau!70!black}{2.2cm},anchor=north west] (a6) at (2.9,-2.6)
    {\jdl{\f{6_topik}}\nodepart{two}sama, + topik $\times$ label};
  \draw[ar] (a1.east)--(a2.west);
  \draw[ar] (a2.east)--(s3.west) node[pos=0.55,above,sloped,font=\tiny] {single};
  \draw[ar] (a2.east)--(m3.west) node[pos=0.55,below,sloped,font=\tiny] {multilabel};
  \draw[ar] (s3.east)--(s4.west); \draw[ar] (s4.east)--(s5.west);
  \draw[ar] (m3.east)--(m4.west); \draw[ar] (m4.east)--(m5.west);
  \draw[ar] (a2.south)--(a6.north);
\end{tikzpicture}
\end{gambar}

\begin{tcolorbox}[enhanced,colback=black!3,colframe=black!25,boxrule=0.5pt,arc=2pt,
  left=8pt,right=8pt,top=5pt,bottom=5pt,fontupper=\small\ttfamily]
\textcolor{jingga}{LABEL\_COLS} = ["harga", "layanan", "kualitas"]\hfill\pil{jingga}{wajib untuk multilabel}\\
\textcolor{hijau}{TARGET} = "harga"\hfill\pil{abu}{satu label terpenting}
\end{tcolorbox}

\begin{itemize}
  \item \textbf{\f{LABEL_COLS}} adalah daftar kolom label 0/1 di train. Selama isinya \f{None},
  notebook \f{m} berhenti dengan pesan galat.
  \item \textbf{\f{TARGET}} tetap dibutuhkan \f{1_eda}, \f{2_praproses}, dan grafik ``topik
  $\times$ kelas'', yang hanya mengerti satu kolom kelas. Isi dengan satu label terpenting; kalau
  tidak diisi, dipakai label pertama.
  \item \textbf{Hasil} masuk \f{baseline_multilabel/}, \f{finetune_multilabel/},
  \f{ensemble_multilabel/}, dan OOF ke \f{oof_multilabel/}. \f{8_ekspor.bat} sudah mengambilnya.
  \item \textbf{\f{6_topik}} menambah grafik \f{06b_topik_x_label.png}: porsi teks berlabel di tiap
  topik beserta lift-nya, plus kalimatnya di \f{fakta_topik.md}.
  \item \textbf{\f{1_eda}} menjalankan bagian \textbf{7f. Multilabel} (dilewati bila \f{LABEL_COLS = None}):
  porsi tiap label, jumlah label per teks, label yang sering muncul bersama (P(kolom | baris)),
  panjang teks per label, teks kembar dengan kombinasi label berbeda. Luaran di \f{<HASIL>/eda/}:
  \f{07f_multilabel.png}, \f{07f_panjang_per_label.csv}, kalimatnya di \f{temuan.md}. Kolom label
  yang bukan 0/1 menghentikan bagian ini dengan pesan galat.
  \item \textbf{Mencoba dulu:} di bagian bawah blok NILAI SOAL tiap notebook ada baris ber-\texttt{\#}
  untuk data demo CASA (6 label). Buka \texttt{\#} di semua notebook yang dijalankan (baris yang
  variabelnya tidak ada di notebook itu tidak berpengaruh).
\end{itemize}

\begin{nb}{kuning}{\f{3m_baseline_multilabel}}{demo CASA 5 s, F1 makro OOF 0,862}
\begin{isi}
Tujuan & Submission multilabel pertama dan OOF pertama, tanpa GPU.\\
Ubah & \f{MODEL_NAME} tiap variasi; \f{MIN_ONE_LABEL = True} bila tiap teks pasti punya minimal
satu label.\\
Lihat & Tabel ambang per label dan \textbf{\f{MASALAH: 0}}. Baris \f{nilai label selain 0/1}
harus 0.
\end{isi}
\end{nb}

\begin{nb}{jingga}{\f{4m_finetune_multilabel}}{demo CASA 150 s (5 epoch, 5 fold), F1 makro OOF 0,924}
\begin{isi}
Tujuan & Skor tertinggi: satu logit per label, loss BCE, ambang per label dari OOF.\\
Ubah & \f{PRETRAINED} + \f{MODEL_NAME}; \f{POS_WEIGHT = True} untuk label langka; \f{EPOCHS}
(bawaan 5, untuk data kecil; data besar cukup 3).\\
Awas & Hanya memakai kolom teks pertama di \f{TEXT_COLS}. Lipatannya sama dengan \f{3m}
(\f{KFold} seed 42), jadi jangan ubah \f{SPLIT_SEED}.
\end{isi}
\end{nb}

\begin{nb}{ungu}{\f{5m_ensemble_multilabel}}{demo CASA 9 s}
\begin{isi}
Tujuan & Menggabung OOF \f{3m} dan \f{4m}, menulis submission akhir, sekaligus laporan untuk
slide.\\
Lihat & Sel ``Putusan''. Rata-rata dipakai hanya bila skor jujur naik sekurangnya \f{MIN_GAIN};
di demo yang dipilih model tunggal \f{ml_indobert} (0,919 lawan rata-rata 0,913).\\
Luaran & \f{submission.csv}, \f{fakta_multilabel.md}, skor per label, jumlah label per teks,
label yang muncul bersama.
\end{isi}
\end{nb}

\begin{awas}[kuning]
\textbf{Label ditulis dalam satu kolom} (misalnya \f{"harga, layanan"})? Notebook \f{m} butuh
kolom 0/1. Pecah dulu di sel baru paling atas \f{2_praproses}, sesudah sel ``Baca data'', lalu isi
\f{LABEL_COLS} dengan nama kolom yang muncul:\par\smallskip
{\ttfamily\small pecah = train["label"].str.replace(" ", "").str.get\_dummies(sep=",")\\
train = pd.concat([train, pecah], axis=1)\\
print(pecah.columns.tolist())}\par\smallskip
Potongan ini tidak ikut diuji dalam rantai; cek hasil \f{print}-nya sebelum lanjut. Format
submission multilabel juga bisa berbeda dari kolom 0/1 (ikuti sample panitia).
\end{awas}

% ================= 4 =================
\section{Jam terakhir: tiga submission}

Empat jam pertama tidak ada leaderboard. Pilih file dari skor OOF dan dari validator, bukan dari
perasaan.

\begin{gambar}{Urutan tiga submission (saran)}
\begin{tikzpicture}
  \node[kartu={kuning}{4.3cm},anchor=north west,font=\footnotesize] (s1) at (0,0)
    {\jdl{Submission 1: yang pasti sah}\nodepart{two}Model tunggal dengan OOF tertinggi.
     Tujuannya memastikan format diterima dan melihat selisih OOF dengan leaderboard.};
  \node[kartu={ungu}{4.3cm},anchor=north west,font=\footnotesize] (s2) at (5.35,0)
    {\jdl{Submission 2: ensemble}\nodepart{two}Keluaran \f{5_ensemble}, bila putusannya memang
     memilih gabungan. Kalau tidak, model tunggal terbaik kedua.};
  \node[kartu={merah}{4.3cm},anchor=north west,font=\footnotesize] (s3) at (10.7,0)
    {\jdl{Submission 3: cadangan}\nodepart{two}Simpan sampai 17.40. Pakai untuk memperbaiki
     kesalahan format, atau untuk kandidat OOF tertinggi berikutnya.};
  \draw[ar] (s1.east)--(s2.west); \draw[ar] (s2.east)--(s3.west);
\end{tikzpicture}
\end{gambar}

\textbf{Sebelum mengunggah tiap file:}\par
\cek validator di notebook penghasilnya menulis \f{MASALAH: 0} dan tidak ada
\f{PERINGATAN: SAMPLE_CSV masih file demo};\par
\cek jumlah baris dan urutan id sama dengan sample panitia (dicek validator);\par
\cek sebaran kelas prediksi masuk akal, tidak satu kelas saja (dicetak validator);\par
\cek nama file dan nama kolom persis seperti yang diminta soal.

Leaderboard publik hanya sebagian data test. Selisih kecil di sana belum tentu nyata; kalau OOF
dan leaderboard tidak sepakat untuk selisih kecil, pegang OOF.

% ================= 5 =================
\section{Yang boleh keluar dari PC lomba}

\begin{minipage}[t]{0.485\linewidth}
\begin{tcolorbox}[enhanced,colback=hijau!6,colframe=hijau,boxrule=0.5pt,arc=2pt,
  title={Boleh dibawa},fonttitle=\bfseries,colbacktitle=hijau,left=6pt,right=6pt,
  equal height group=nda,fontupper=\small\raggedright]
\begin{itemize}
  \item PNG grafik dan tabel ringkas
  \item \f{temuan.md}, \f{fakta_topik.md}, \f{fakta_model.md}
  \item notebook, file submission, model
  \item screenshot
\end{itemize}
Semuanya dikumpulkan \f{8_ekspor.bat}.
\end{tcolorbox}
\end{minipage}\hfill
\begin{minipage}[t]{0.485\linewidth}
\begin{tcolorbox}[enhanced,colback=merah!6,colframe=merah,boxrule=0.5pt,arc=2pt,
  title={Dilarang keluar},fonttitle=\bfseries,colbacktitle=merah,left=6pt,right=6pt,
  equal height group=nda,fontupper=\small\raggedright]
\begin{itemize}
  \item teks asli atau baris data
  \item \f{train_bersih.csv}, \f{test_bersih.csv}
  \item file per baris: OOF (\f{.npy}), \f{topik_dokumen.csv}, \f{pseudo_label.csv}
  \item output sel yang mencetak teks
\end{itemize}
\end{tcolorbox}
\end{minipage}

\smallskip
Slide dibuat di laptop tanpa akses data, jadi setiap angka yang ingin disebut harus sudah ada di
tiga file \f{.md} di atas atau di PNG. Kalau ada angka yang kurang, tambahkan ke notebook
\emph{sebelum} meninggalkan PC lomba.

% ================= 6 =================
\section{Kalau macet}

{\small
\begin{tabularx}{\linewidth}{@{}>{\raggedright\arraybackslash}p{5.2cm}Y@{}}
\toprule
\textbf{Gejala} & \textbf{Tindakan}\\
\midrule
\f{CUDA out of memory} & \f{BATCH_SIZE = 8} dan \f{ACCUM_STEPS = 2}; kecilkan \f{MAX_LENGTH};
tutup kernel notebook lain yang memakai GPU (termasuk \f{6_topik}).\\
\f{device: cpu} di \f{4_finetune} & Driver NVIDIA belum cocok (torch cu128 butuh driver 570 ke
atas). Selagi diperbaiki, kerjakan variasi \f{3_baseline}.\\
Satu fold terlalu lama & \f{distilbert_id} atau \f{minilm}; \f{EPOCHS = 2}; \f{PATIENCE = 2}
dengan \f{EVALS_PER_EPOCH = 2}.\\
\f{TIDAK ADA OOF YANG COCOK} & OOF dibuat dari train yang lain: \f{2_praproses} dijalankan ulang
atau \f{HASIL} masih memuat OOF demo. Latih ulang model, atau ganti nama \f{HASIL}.\\
\f{id test = id sample: False} & \f{ID_COL} salah, atau test bukan file yang dipasangkan dengan
sample.\\
\f{UnicodeDecodeError}, kolom jadi satu & Baris ber-\texttt{\#} di sel ``Baca data'' \f{2_praproses}
(\f{encoding}, \f{sep}).\\
Prediksi hampir satu kelas & Kelas timpang: \f{CLASS_WEIGHT = True} di \f{4_finetune}; baseline
sudah memakai \f{class_weight="balanced"} dan pengali kelas.\\
OOF tinggi, leaderboard rendah & Curigai teks kembar train (isi \f{GROUP_COL} atau lihat temuan
kembar di \f{1_eda}) dan pergeseran train--test (AUC adversarial di \f{1_eda}).\\
Topik berisi kata sampah & Di sel ``Konfigurasi'' \f{6_topik.ipynb}: tambah ke
\f{EXTRA_STOPWORDS}, naikkan \f{MIN_DF}, turunkan \f{MAX_DF}. Kata dasar aneh (stemming salah,
mis.\ nama tempat terpotong): \f{TOPIK_BERAT = False} di blok NILAI SOAL \f{6_topik.ipynb}, lalu
Run All \f{6_topik}.\\
\f{KeyError: 'text_tfidf'} di \f{6_topik} & \f{<HASIL>/praproses/train_bersih.csv} tidak punya
kolom \f{<teks>_tfidf} (file lama atau \f{TEXT_COLS} berubah): Run All \f{2_praproses.ipynb},
lalu \f{6_topik.ipynb}.\\
\bottomrule
\end{tabularx}}

% ================= 7 =================
\section{Daftar cek}

\begin{minipage}[t]{0.485\linewidth}\small
\textbf{Sebelum 13.00}\par
\cek susunan folder seperti Gambar 3\par
\cek \f{jupyter.bat} membuka JupyterLab\par
\cek rantai 1--7 pernah Run All di data demo\par
\cek \f{4_finetune} menulis \f{device: cuda}\par
\cek folder \f{outputs/final} demo tidak dipakai lagi\par
\end{minipage}\hfill
\begin{minipage}[t]{0.485\linewidth}\small
\textbf{Sebelum meninggalkan PC}\par
\cek tiga file submission tercatat: nama, skor OOF, skor leaderboard\par
\cek \f{6_topik}: \f{TOPIC_NAMES} sudah diisi\par
\cek \f{7_laporan_model} dijalankan sesudah model terakhir\par
\cek \f{8_ekspor.bat}, lalu baca \f{MANIFEST.md}\par
\cek screenshot leaderboard\par
\end{minipage}

\end{document}
